\begin{abstract}
Model merging in weight space is a powerful paradigm for combining multiple task-specific expert models into a single unified network without incurring additional training costs. However, standard linear interpolation (Task Arithmetic) and sign-consistent average merging (TIES-Merging) inevitably suffer from spatial weight-space interference, causing severe representational collapse when merging highly orthogonal task models. To mitigate this, recent literature has explored test-time adaptation (TTA) frameworks, optimization networks, and continuous layer-wise scaling coefficients. While powerful, these approaches often introduce highly parameterized, multi-stage optimization pipelines. In this work, we show that weight-space interference can be addressed through a direct coordinate-level allocation strategy. We propose \textbf{Exclusive Parameter Merging (EPM)}, a training-free model fusion operator that mitigates weight-space interference by introducing a soft coordinate-wise assignment strategy. EPM introduces \textbf{Soft Exclusive Parameter Allocation (Soft-EPA)}, which routes coordinate updates primarily to the dominant task expert based on scaled absolute magnitude, while attenuating non-dominant updates by a coherence retention factor $\gamma = 0.2$. This balances coordinate exclusivity and cooperative blending, suppressing weight-space conflicts while preserving structural representation coherence across layers. To optimize the task-specific scales while avoiding the overfitting-optimizer paradox of high-dimensional tuning, we introduce \textbf{Task-Level Coefficient Tuning (TLC-Tune)}, which optimizes only $K$ global scaling factors via a stable, balanced multi-task zero-order (1+1) Evolution Strategy on a modest offline validation split of 128 samples per task. Empirically, EPM and TLC-Tune demonstrate outstanding multi-task performance and resilience on a Vision Transformer (ViT-Tiny) backbone across four conflicting image datasets (MNIST, FashionMNIST, CIFAR-10, SVHN). We compare EPM against traditional averaging methods as well as highly parameterized layer-wise optimization frameworks (AdaMerging and ZipMerge). EPM consistently outclasses all baselines at 0\% and 50\% sparsities (achieving up to \textbf{46.19\%} joint mean accuracy), while empirically validating the Overfitting-Optimizer Paradox by showing that higher-dimensional validation-set tuning severely degrades final test generalization.
\end{abstract}
